\documentclass[paper=a4, fontsize=11pt]{jlreq} % 日本語用の標準的なクラス

% --- パッケージ群 ---
\usepackage[utf8]{inputenc}
\usepackage[T1]{fontenc}
\usepackage{amsmath, amssymb, amsthm} % 数学の基本
\usepackage{physics} % \ket{0}, \pdv{x} などが簡単に打てる
\usepackage{bm} % 太字ベクトル \bm{v}
\usepackage{graphicx}
\usepackage{hyperref} % リンク・目次用
\usepackage{cite} % 引用用

% --- 定理環境の設定 (OMCの問題風に整理できる) ---
\theoremstyle{definition}
\newtheorem{theorem}{定理}[section]
\newtheorem{definition}[theorem]{定義}
\newtheorem{example}[theorem]{例}
\newtheorem{remark}[theorem]{注釈}
\newtheorem{problem}[theorem]{問題} % 自作問題用

% --- 独自のショートカット (CFT/パンルヴェ用) ---
\newcommand{\Virexp}[2]{[L_{#1}, L_{#2}]} % ヴィラソロの交換関係
\newcommand{\taufunc}{\tau(z)} % タウ関数

\title{多様体に関する勉強ノート}
\author{松浦光佑}
\date{\today}

\begin{document}

\maketitle
\tableofcontents % 自動で目次を作成

\section{はじめに}
本ノートは、松本『多様体の基礎』に関する学習記録である。
日々の学習を「AC（Accepted）」として記録し、GitHubでの進捗可視化を目的とする。

\section{多様体の基礎}
\begin{definition}[座標近傍、局所座標系]
    % --- ここに数式 ---
    X:位相空間とする。$U\subset X$:開集合から$U'\subset\mathbb{R}^{m}$:開集合への同相写像
    \[
    \phi : U \to U'
    \]
    があるとき$(U,\phi)$を\textbf{m次元座標近傍}といい、$\phi$を$U$上の\textbf{局所座標系}という。
    \begin{itemize}
        \item[\textbf{Insight:}] 多様体上の小領域を数空間$\mathbb{R}^{m}$に対応させる「地図」であり、抽象的な図形の上で微分などの計算を可能にするための定義。地球儀の一部を平面図として扱うように、図形上の各点に「座標」という数値の番地を割り振る役割を担っている。
        \item[\textbf{Link:}] 松本6節の定義6.1/p.37を参照。
    \end{itemize}
\end{definition}

\begin{definition}[位相多様体]
    % --- ここに数式 ---
    位相空間$M$が次の条件(1)(2)を満たす時、$M$を\textbf{m次元位相多様体}という。
    \begin{itemize}
        \item[(1)] $M$はハウスドルフ空間
        \item[(2)] 任意の$p\in M$に対して$p$を含むm次元座標近傍$(U,\phi)$が存在する。
    \end{itemize}
    \begin{itemize}
        \item[\textbf{Insight:}] 「局所的には数空間 $\mathbb{R}^m$ と見なせる」という性質により、複雑な図形を慣れ親しんだ計算の土俵に載せるための定義。ハウスドルフ性は空間としての健全性（点が適切に分離できること）を保証し、全域が「地図（座標近傍）」で漏れなく覆い尽くせることを約束している。
        \item[\textbf{Link:}] 松本6節の定義6.2/p.38を参照。
    \end{itemize}
\end{definition}

\begin{definition}[座標変換]
    % --- ここに数式 ---
    $(U,\varphi),(V,\psi)$に対して、同相写像$\psi\circ\varphi^{-1}:\varphi(U\cap V)\to\psi(U\cap V)$を、$(U,\varphi)$から$(V,\psi)$への\textbf{座標変換}と呼ぶ
    \begin{itemize}
        \item[\textbf{Type:}] $(U,\varphi),(V,\psi)$はm次元多様体$M$の座標近傍
        \item[\textbf{Insight:}] 座標変換 $\psi \circ \varphi^{-1}$ は、多様体の重なり合う2つの「地図」の間で、座標値を直接翻訳するルールである。これにより、抽象的な曲がった空間上の問題を、私たちが扱い慣れたユークリッド空間同士の具体的な計算へと落とし込むことができる。
        \item[\textbf{Link:}] 松本6節の定義6.3/p.41を参照。
    \end{itemize}
\end{definition}

\begin{example}[直交座標と極座標の座標変換]
    平面 $\mathbb{R}^2$ の第1象限を開集合 $U$ とする。$U$ 上の点に対して、2組の局所座標系として直交座標 $(x,y)$ と極座標 $(r,\theta)$ を考える。
    
    極座標を与える座標近傍を $(U, \varphi)$、直交座標を与える座標近傍を $(U, \psi)$ としたとき、$(U, \varphi)$ から $(U, \psi)$ への座標変換 $\psi \circ \varphi^{-1} : (r, \theta) \mapsto (x, y)$ は以下の式で与えられる。
    \begin{align*}
        x &= r \cos \theta \\
        y &= r \sin \theta
    \end{align*}
    
    \begin{itemize}
        \item[\textbf{Insight:}] 抽象的な記号 $\psi \circ \varphi^{-1}$ の正体は、微積分でよく知っている「極座標から直交座標への変換公式」そのものである。多様体論における座標変換が、決して特殊な操作ではなく、馴染み深い概念を厳密に定式化し直したものに過ぎないことがよくわかる。
        \item[\textbf{Link:}] 松本6節の例/p.41を参照。
    \end{itemize}
\end{example}

\begin{definition}[$C^{r}$級微分可能多様体]
    % --- ここに数式 ---
    $r\geq 1$を自然数または$\infty$とする。位相空間$M$が次の条件(1)-(3)を満たす時、$M$を\textbf{m次元$C^{r}$級微分可能多様体}という。\textbf{m次元$C^{r}$級多様体}ともいう。
    \begin{itemize}
        \item[(1)] $M$はハウスドルフ空間。
        \item[(2)] $M$はm次元の座標近傍によって被覆される。すなわち、$M$はm次元座標近傍からなる族$\{(U_{\alpha},\varphi_{\alpha})\}_{\alpha\in A}$があって、
                    \[
                    M = \bigcup_{\alpha\in A}U_{\alpha}
                    \]
                    が成り立つ。($A$は適当な添字集合)\\
                    $\{(U_{\alpha},\varphi_{\alpha})\}_{\alpha\in A}$を\textbf{座標近傍系}または\textbf{アトラス}という。
        \item[(3)] $U_{\alpha}\cap U_{\beta}\neq\phi$であるような任意の$\alpha,\beta\in A$について、座標変換
                    \[
                    \varphi_{\beta}\circ\varphi_{\alpha}^{-1}:\varphi_{\alpha}(U_{\alpha}\cap U_{\beta})\to\varphi_{\beta}(U_{\alpha}\cap U_{\beta})
                    \]
                    は$C^{r}$級写像である。\\
                    この座標近傍系$\{(U_{\alpha},\varphi_{\alpha})\}_{\alpha\in A}$を\textbf{$C^{r}$級座標近傍系}という。
    \end{itemize}
    \begin{itemize}
        \item[\textbf{Insight:}] ただの「つぎはぎの地図（位相多様体）」を、微積分が可能な「滑らかな空間」へとアップグレードする定義。地図同士のつなぎ目（座標変換）が $C^r$ 級で滑らかに接続されることを要求することで、多様体全体で矛盾なく微分や接ベクトルの計算ができるようになる。
        \item[\textbf{Link:}] 松本6節の定義6.4/p.42を参照。
    \end{itemize}
\end{definition}

\begin{example}[多様体の例$(\mathbb{R}^{m})$]
    $\mathbb{R}^{m}$は$m$次元多様体である。
    $(U, \phi) = (\mathbb{R}^{m}, id)$とすればよい。
    これが定義を満たすことは以下のように確認できる：
    \begin{itemize}
        \item[(1)] $\mathbb{R}^m$はユークリッド空間の位相によりハウスドルフ空間である。
        \item[(2)] $\mathbb{R}^m$自身が開集合であり、$\mathrm{id}: \mathbb{R}^m \to \mathbb{R}^m$ は同相写像。座標近傍系 $\{(U, \phi)\}$ は明らかに $M=\mathbb{R}^m$ 全体を被覆する。
        \item[(3)] 座標近傍が1つしかないため、座標変換は $\mathrm{id} \circ \mathrm{id}^{-1} = \mathrm{id}$ となる。これはすべての成分が1次式なので、明らかに $C^\infty$ 級（したがって $C^r$ 級）である。
    \end{itemize}   
    \begin{itemize}
        \item[\textbf{Insight:}] 多様体論が手本としている「基準となる空間」そのものであり、最も自明な例。空間全体をたった1枚の「地図（恒等写像 $id$）」で歪みなく覆いつくせるため、つなぎ目（座標変換）を考える必要すらない。多様体の複雑な定義が、我々のよく知る平坦なユークリッド空間を矛盾なく包摂している（自然な拡張になっている）ことを保証する、いわば理論のベースケースである。
        \item[\textbf{Link:}] 松本6節の例1/p.43を参照。
    \end{itemize}
\end{example}




% --- 今後の課題（TODOリスト） ---
\section{Next Step / To-Do}
\begin{itemize}
    \item [ ] 多様体の例を見ていく
\end{itemize}

\end{document}